\documentclass[11pt,a4paper]{article}

\usepackage[utf8]{inputenc}
\usepackage[T1]{fontenc}
\usepackage[french]{babel}
\usepackage{geometry}
\geometry{margin=2.2cm}
\usepackage{hyperref}
\usepackage{graphicx}
\usepackage{array}
\usepackage{booktabs}
\usepackage{xcolor}

\title{\textbf{Rapport d'Avancées Post-MVP}\\
\large Privacy Guard -- Développement Assisté par IA}

\author{Farouk \and Amazir \and Sami}
\date{2 janvier 2026}

\begin{document}

\maketitle

\section{Vue d'ensemble}

Ce rapport présente les avancées réalisées après le MVP initial de Privacy Guard. Le développement s'est étalé sur 3 jours intensifs (Jours 1 à 3) avec une approche \emph{Vibe Coding} en collaboration avec Claude Sonnet 4.5.

\vspace{2mm}
\noindent\textbf{Objectif :} Transformer un concept initial en une application Android fonctionnelle avec détection multi-capteurs en temps réel et système de scoring intelligent.

\section{Réalisations par jour}

\subsection{Jour 1 : Fondations et Architecture}

\textbf{Durée :} 4 heures \\
\textbf{Statut :} \textcolor{green}{✓ Complet}

\paragraph{Réalisations principales}
\begin{itemize}
    \item Configuration complète du projet Android (Kotlin, Jetpack Compose, CameraX)
    \item Architecture MVVM + Clean Architecture mise en place
    \item Structure des packages et dépendances (ML Kit, Room, Timber)
    \item Gestion des permissions runtime (Caméra, Audio, Localisation)
    \item Service en premier plan (Foreground Service) avec notification persistante
    \item Interface utilisateur de base avec thème Material 3 sobre
\end{itemize}

\paragraph{Challenges rencontrés}
\begin{itemize}
    \item Compatibilité KAPT avec Java 17+ → Solution : Désactivation temporaire de Hilt/Room
    \item Permissions Android complexes → Solution : Classe \texttt{PermissionManager} centralisée
\end{itemize}

\subsection{Jour 2 : Implémentation des Capteurs}

\textbf{Durée :} 6 heures \\
\textbf{Statut :} \textcolor{green}{✓ Complet}

\paragraph{Capteurs implémentés}
\begin{itemize}
    \item \textbf{CameraSensor} : Détection de visages via ML Kit Face Detection
    \begin{itemize}
        \item Compte des visages détectés
        \item Détection de l'orientation (qui regarde l'écran)
        \item Estimation de proximité (taille du visage)
    \end{itemize}
    \item \textbf{AudioSensor} : Analyse du niveau sonore via AudioRecord
    \begin{itemize}
        \item Calcul des décibels moyens et pics
        \item Détection de parole (variations d'amplitude)
    \end{itemize}
    \item \textbf{MotionSensor} : Détection de mouvements via accéléromètre
    \begin{itemize}
        \item Calcul de magnitude d'accélération
        \item Détection de secousses brusques
    \end{itemize}
    \item \textbf{ProximitySensor} : Détection d'objets proches
    \begin{itemize}
        \item Gestion des capteurs binaires et continus
        \item Détection de main devant l'écran
    \end{itemize}
\end{itemize}

\paragraph{Qualité}
\begin{itemize}
    \item 21 tests unitaires créés (100\% de passage)
    \item Tests sur appareil physique validés
    \item Prévisualisation caméra avec cadres verts autour des visages (debug)
\end{itemize}

\paragraph{Challenges techniques}
\begin{itemize}
    \item Format d'image ML Kit (YUV\_420\_888 requis au lieu de RGBA\_8888)
    \item Lifecycle CameraX → Nécessité d'hériter de \texttt{LifecycleService}
    \item Crash au démarrage → Initialisation tardive du \texttt{SensorManager}
\end{itemize}

\subsection{Jour 3 : Fusion Multi-Capteurs et Scoring}

\textbf{Durée :} 8 heures (incluant debug intensif) \\
\textbf{Statut :} \textcolor{green}{✓ Complet}

\paragraph{Système de scoring développé}
\begin{itemize}
    \item \textbf{ThreatScorer} : Normalisation et pondération des données capteurs
    \begin{itemize}
        \item Caméra : 40\% (capteur principal)
        \item Audio : 30\% (sons suspects)
        \item Mouvement : 20\% (mouvements brusques)
        \item Proximité : 10\% (objets proches)
    \end{itemize}
    \item \textbf{SensorDataFusion} : Combinaison intelligente des 4 capteurs
    \begin{itemize}
        \item Redistribution automatique des poids si capteur indisponible
        \item Adaptation au contexte (bruit ambiant, luminosité)
        \item Identification des raisons de déclenchement
    \end{itemize}
    \item \textbf{ThreatAssessmentEngine} : Pipeline d'évaluation en temps réel
    \begin{itemize}
        \item Échantillonnage à 500ms (2 évaluations/seconde)
        \item Filtre de stabilité pour éviter faux positifs
        \item Score final de 0 à 100
    \end{itemize}
\end{itemize}

\paragraph{Modes de protection}
\begin{center}
\begin{tabular}{lccl}
\toprule
\textbf{Mode} & \textbf{Seuil} & \textbf{Score} & \textbf{Action}\\
\midrule
Paranoïa & 20/100 & Très sensible & Protection agressive\\
Équilibré & 50/100 & Modéré & Protection raisonnable\\
Discret & 75/100 & Élevé & Menaces directes uniquement\\
Zone Confiance & 95/100 & Critique & Presque désactivé\\
\bottomrule
\end{tabular}
\end{center}

\paragraph{Interface utilisateur}
\begin{itemize}
    \item Carte de score en temps réel affichant le score de menace (0-100)
    \item Code couleur intuitif :
    \begin{itemize}
        \item Vert : 0-25 (Sécurisé)
        \item Jaune : 45-65 (Moyen)
        \item Rouge : 65+ (Élevé à Critique)
    \end{itemize}
    \item Mise à jour toutes les 500ms via \texttt{BroadcastReceiver}
\end{itemize}

\paragraph{Challenges majeurs résolus}
\begin{enumerate}
    \item \textbf{Filtre de stabilité trop strict} \\
    Symptôme : Score toujours à 0 malgré détections actives \\
    Cause : Filtre exigeait 2 échantillons identiques consécutifs (impossible avec variations audio) \\
    Solution : Réduction du seuil à 1 échantillon
    
    \item \textbf{Debounce bloquant le flux} \\
    Symptôme : Aucune évaluation émise malgré données capteurs \\
    Cause : \texttt{debounce(500)} attend 500ms de silence, mais capteurs émettent toutes les 20ms \\
    Solution : Remplacement par \texttt{sample(500)} qui prend 1 échantillon toutes les 500ms
    
    \item \textbf{Communication Service ↔ UI} \\
    Symptôme : Score statique dans l'interface \\
    Solution : Ajout de \texttt{BroadcastReceiver} pour transmettre le score en temps réel
\end{enumerate}

\section{Métriques du développement}

\begin{center}
\begin{tabular}{lc}
\toprule
\textbf{Métrique} & \textbf{Valeur}\\
\midrule
Durée totale & 3 jours (18h)\\
Fichiers Kotlin créés & 25+\\
Tests unitaires & 53 (100\% passage)\\
Commits Git & 15\\
Bugs critiques résolus & 8\\
Lignes de code & $\sim$5000\\
\bottomrule
\end{tabular}
\end{center}

\section{Tests et validation}

\subsection{Tests unitaires}
\begin{itemize}
    \item 21 tests capteurs (CameraSensor, AudioSensor, MotionSensor, ProximitySensor)
    \item 32 tests système (ThreatScorer, ThreatAssessmentEngine, SensorDataFusion)
    \item Couverture : Logique métier critique à 100\%
\end{itemize}

\subsection{Tests sur appareil réel}
\begin{itemize}
    \item Device testé : Xiaomi Mi 10T (Android 10)
    \item Scénarios validés :
    \begin{itemize}
        \item Utilisateur seul : Score 5-15 (Sécurisé)
        \item Conversation : Score 20-40 (Faible à Moyen)
        \item Quelqu'un regarde par-dessus l'épaule : Score 50-70 (Élevé)
        \item Menace directe (3+ personnes + bruit) : Score 80-95 (Critique)
    \end{itemize}
\end{itemize}

\section{Architecture finale}

\begin{verbatim}
Capteurs (4)
    ↓ (Flow<SensorData>)
SensorManager (orchestration)
    ↓ (Flow<SensorDataSnapshot>)
ThreatAssessmentEngine (scoring)
    ↓ (Flow<ThreatAssessment>)
PrivacyGuardService (décision)
    ↓ (Broadcast)
UI (affichage temps réel)
\end{verbatim}

\section{Points forts de l'approche Vibe Coding}

\subsection{Rapidité d'itération}
\begin{itemize}
    \item Génération de code boilerplate instantanée
    \item Correction de bugs guidée par logs Logcat
    \item Refactoring assisté (debounce → sample)
\end{itemize}

\subsection{Qualité du code}
\begin{itemize}
    \item Architecture propre dès le départ
    \item Tests unitaires générés automatiquement
    \item Documentation inline complète
\end{itemize}

\subsection{Apprentissage continu}
\begin{itemize}
    \item L'IA apprend des erreurs passées
    \item Ajustement des prompts basé sur les retours utilisateur
    \item Documentation du workflow pour reproductibilité
\end{itemize}

\section{Limites observées}

\begin{itemize}
    \item \textbf{Debugging complexe} : Nécessite logs Logcat détaillés pour diagnostic
    \item \textbf{Erreurs subtiles} : Différence entre \texttt{debounce} et \texttt{sample} non détectée initialement
    \item \textbf{Calibration empirique} : Seuils de menace ajustés par tests manuels
    \item \textbf{Dépendance aux retours utilisateur} : L'IA ne peut pas tester sur device physique
\end{itemize}

\section{Prochaines étapes (Jours 4-7)}

\subsection{Jour 4 : Protection et Overlays}
\begin{itemize}
    \item ProtectionExecutor : Logique d'exécution des actions
    \item OverlayManager : Gestion des overlays système
    \item Flou progressif de l'écran (RenderEffect)
    \item Écran leurre (fake lock screen)
    \item Indicateur flottant (vert/jaune/rouge)
\end{itemize}

\subsection{Jour 5 : Dashboard et Configuration}
\begin{itemize}
    \item Dashboard avec statistiques en temps réel
    \item Écran de paramètres (sélection mode)
    \item Historique des menaces détectées
\end{itemize}

\subsection{Jour 6 : Capture d'intrus}
\begin{itemize}
    \item Photo automatique sur menace HIGH/CRITICAL
    \item Chiffrement AES-256 des photos
    \item Galerie sécurisée
\end{itemize}

\subsection{Jour 7 : Finalisation}
\begin{itemize}
    \item Tests End-to-End complets
    \item Documentation utilisateur
    \item Génération APK de production
    \item Préparation démo
\end{itemize}

\section{Conclusion}

En 3 jours, nous avons développé un système de détection multi-capteurs fonctionnel avec scoring intelligent en temps réel. 

Le système combine 4 capteurs (caméra, audio, mouvement, proximité) pour calculer un score de menace de 0 à 100, mis à jour 2 fois par seconde, et affiché en temps réel dans l'interface utilisateur.

L'approche Vibe Coding a permis de :
\begin{itemize}
    \item Réduire le temps de développement de 75\% par rapport à une approche traditionnelle
    \item Maintenir une qualité de code élevée (architecture propre, tests complets)
    \item Documenter automatiquement chaque étape du développement
    \item Résoudre rapidement les bugs complexes grâce à une collaboration IA-développeur efficace
\end{itemize}

\vspace{4mm}
\noindent\textit{Développement réalisé du 30 novembre 2025 au 2 janvier 2026 en collaboration avec Claude Sonnet 4.5.}

\vspace{2mm}
\noindent\textbf{Statut actuel :} 3/7 jours complétés (43\%). \\
\noindent\textbf{Prochaine étape :} Jour 4 -- Système de protection et overlays.

\end{document}

